%% ════════════════════════════════════════════════════════════════════════
%%  E₁ MODULAR KOSZUL DUALITY
%%
%%  The ordered (associative) parallel to the commutative modular
%%  Koszul theory.  Extracted to a standalone chapter so that the
%%  E₁/E_∞ parallelism is architecturally visible.
%% ════════════════════════════════════════════════════════════════════════

\chapter{\texorpdfstring{$E_1$}{E1} Modular Koszul Duality}
\label{chap:e1-modular-koszul}

Replacing the commutative modular operad $\cM_{\Com}$ and its Feynman
transform $F\!\Com$ by the ribbon modular operad $\cM_{\Ass}$ and
$F\!\Ass$ yields the associative ($E_1$) modular Koszul theory.

The structural parallel is exact.  Every construction in the
commutative theory (the modular convolution algebra $\gAmod$, the
bar-intrinsic MC element $\Theta_\cA$, the shadow Postnikov tower,
and the five main theorems A--H) has an ordered counterpart governed
by $F\!\Ass$ in place of $F\!\Com$.  The $E_\infty$ invariants are
recovered as $\Sigma_n$-coinvariants of the $E_1$ invariants via the
averaging map
(Definition~\ref{def:e1-modular-convolution}).
The additional data visible at the $E_1$ level (the spectral
$R$-matrix $r(z)$, the Drinfeld associator as the arity-$3$ shadow
of affine Kac--Moody, the ordered genus-$1$ monodromy) are genuine
refinements invisible after averaging.

\medskip
\noindent\textbf{Plan.}
Section~\ref{sec:e1-associative-modular-operad} constructs the
ribbon modular operad, its Feynman transform, and the $E_1$ modular
convolution algebra.
Section~\ref{sec:e1-shadow-tower} develops the $E_1$ shadow
Postnikov tower and identifies the low-arity shadows (CYBE,
pentagon, KZ associator).
Section~\ref{sec:e1-five-theorems} proves the five main theorems
$\mathrm{A}^{E_1}$--$\mathrm{H}^{E_1}$ at all genera.


%% ════════════════════════════════════════════════════════════════════════
%%  THE ASSOCIATIVE MODULAR OPERAD AND E_1 CONVOLUTION
%% ════════════════════════════════════════════════════════════════════════

\section{The associative modular operad and \texorpdfstring{$E_1$}{E1}
convolution}
\label{sec:e1-associative-modular-operad}
\index{associative modular operad|textbf}
\index{FAss@$F\Ass$|textbf}
\index{ribbon modular operad|textbf}
\index{E1 algebra@$E_1$-algebra!modular convolution}

\begin{definition}[Ribbon modular operad {\cite{GeK98,CG17}}]
\label{def:ribbon-modular-operad}
\ClaimStatusProvedElsewhere
\index{ribbon modular operad!definition}
The \emph{ribbon modular operad}
$\cM_{\Ass}$ has components
\begin{equation}
\label{eq:ribbon-modular-operad}
\cM_{\Ass}(g,n)
\;:=\;
C_*\!\bigl(\overline{\cM}_{g,n}^{\,\mathrm{rib}}\bigr),
\qquad 2g - 2 + n > 0,
\end{equation}
where $\overline{\cM}_{g,n}^{\,\mathrm{rib}}$ is the moduli space
of stable genus-$g$ Riemann surfaces with~$n$ parametrised boundary
circles.  A ribbon structure on a stable graph $\Gamma$ is a cyclic
ordering of the half-edges at each vertex; equivalently,
$\overline{\cM}_{g,n}^{\,\mathrm{rib}}$ classifies surfaces-with-%
pants-decomposition in which each pair-of-pants carries a boundary
parametrisation.  The composition maps are:
\begin{enumerate}[label=\textup{(\roman*)}]
\item \emph{Self-sewing} $\xi_{ij}\colon
  \cM_\Ass(g,n) \to \cM_\Ass(g+1,n-2)$: identifies the
  $i$-th and $j$-th boundary circles with reversed orientation.
\item \emph{Grafting} $\circ_i\colon
  \cM_\Ass(g_1,n_1) \otimes \cM_\Ass(g_2,n_2)
  \to \cM_\Ass(g_1+g_2, n_1+n_2-2)$: sews the $i$-th
  boundary of the first surface to a boundary of the second.
\end{enumerate}
In genus~$0$: $\cM_{\Ass}(0,n) = C_*(K_n)$, where $K_n$ is the
Stasheff associahedron (dimension~$n-2$), and grafting is face
inclusion.  The forgetful map that remembers the surface but forgets the
ordering of the $n$ external boundary components gives a
$\Sigma_n$-quotient
$\overline{\cM}_{g,n}^{\,\mathrm{rib}} \to
\overline{\cM}_{g,n}^{\,\mathrm{rib}}/\Sigma_n$;
forgetting the ribbon structure at vertices (cyclic orderings of
half-edges) is a separate operation.
\end{definition}

\begin{definition}[Feynman transform of the associative modular operad]
\label{def:feynman-transform-ass}
\ClaimStatusProvedHere
\index{FAss@$F\Ass$!definition}
The \emph{associative Feynman transform} $F\!\Ass$ is the modular
cooperad with $(g,n)$-component
\begin{equation}
\label{eq:fass-definition}
(F\!\Ass)(g,n)
\;:=\;
\bigoplus_{\Gamma \in \mathsf{Gr}^{\mathrm{rib}}_{g,n}}
\frac{1}{|\!\operatorname{Aut}(\Gamma)|}\,
\det\bigl(E(\Gamma)\bigr)
\;\otimes\;
\bigotimes_{v \in V(\Gamma)}
\cM_{\Ass}\bigl(g(v), n(v)\bigr)^{\vee},
\end{equation}
summed over connected ribbon graphs $\Gamma$ of genus~$g$ with~$n$
legs.  The differential
\begin{equation}
\label{eq:fass-differential}
D_{F\!\Ass}
\;=\;
\sum_{e \in E(\Gamma)}
\pm\; \xi_e^*
\end{equation}
is dual to edge contraction in $\cM_\Ass$.
\end{definition}

\begin{theorem}[\ClaimStatusProvedHere]
\label{thm:fass-d-squared-zero}
\index{FAss@$F\Ass$!$D^2=0$}
$D_{F\!\Ass}^2 = 0$.
\end{theorem}

\begin{proof}
$D^2$ pairs with codimension-$2$ boundary strata of
$\overline{\cM}_{g,n}^{\,\mathrm{rib}}$.  Each
codimension-$2$ stratum $S$ is reached by two codimension-$1$
paths (contracting edge~$e_1$ then~$e_2$, or~$e_2$ then~$e_1$),
and the edge-orientation determinant $\det(E)$ supplies a minus
sign between them.  The cancellation uses only the associativity
of face composition in the face poset
(Getzler--Kapranov \cite{GK98}), which holds for ribbon graphs
by the same argument as for ordinary stable graphs: ribbon
structure does not alter the codimension-$2$ face pairing.
\end{proof}

\begin{theorem}[\ClaimStatusProvedHere]
\label{thm:fcom-coinvariant-fass}
\index{FAss@$F\Ass$!abelianization}
$(F\!\Ass)(g,n) / \Sigma_n
\;\xrightarrow{\;\sim\;}
(F\!\Com)(g,n)$.
\end{theorem}

\begin{proof}
The ribbon graph category $\mathsf{Gr}^{\mathrm{rib}}_{g,n}$ is a
principal $\Sigma_n$-bundle over
$\mathsf{Gr}^{\mathrm{st}}_{g,n}$: forgetting the cyclic ordering
at each vertex gives the quotient.  Edge contraction is
$\Sigma_n$-equivariant, so $D_{F\!\Ass}$ descends.
\end{proof}

\begin{definition}[$E_1$ modular convolution dg~Lie algebra]
\label{def:e1-modular-convolution}
\ClaimStatusProvedHere
\index{E1 algebra@$E_1$-algebra!modular convolution|textbf}
\index{modular convolution dg Lie algebra!E1@$E_1$ variant}
Let $\cA$ be a cyclic $E_1$-chiral algebra
(Definition~\ref{def:e1-chiral}).
\begin{equation}
\label{eq:e1-modular-convolution}
{\gAmod}^{E_1}
\;:=\;
\prod_{\substack{g,n \\ 2g-2+n > 0}}
\operatorname{Hom}\!\bigl(
\cM_{\Ass}(g,n),\,
\operatorname{End}_{\cA}(n)
\bigr).
\end{equation}
The Hom is without $\Sigma_n$-equivariance.
The dg~Lie structure is inherited from $F\!\Ass$: $D$ from
$D_{F\!\Ass}$, $[-,-]$ from ribbon-graph composition.
The \emph{averaging map}
\begin{equation}
\label{eq:e1-to-einfty-projection}
\operatorname{av}\colon
{\gAmod}^{E_1}
\;\twoheadrightarrow\;
\gAmod,
\qquad
\operatorname{av}(\phi)(g,n)
\;:=\;
\frac{1}{n!}\sum_{\sigma \in \Sigma_n}
\sigma \cdot \phi(g,n)
\end{equation}
is a surjective dg~Lie morphism.
\end{definition}

\begin{theorem}[$E_1$ Maurer--Cartan element; \ClaimStatusProvedHere]
\label{rem:e1-mc-element}
\index{E1 algebra@$E_1$-algebra!Maurer--Cartan element}
The element
\begin{equation}
\label{eq:theta-e1}
\Theta_\cA^{E_1}
\;:=\;
D_\cA^{E_1} - \dzero
\;\in\;
{\gAmod}^{E_1}
\end{equation}
satisfies the Maurer--Cartan equation
$D\Theta^{E_1} + \tfrac{1}{2}[\Theta^{E_1}, \Theta^{E_1}] = 0$.
Under averaging\textup{:}
\begin{equation}
\label{eq:e1-to-einfty-mc}
\operatorname{av}\bigl(\Theta_\cA^{E_1}\bigr)
\;=\;
\Theta_\cA.
\end{equation}
\end{theorem}

\begin{proof}
$D_{F\!\Ass}^2 = 0$
(Theorem~\ref{thm:fass-d-squared-zero}) and the bar-intrinsic
construction (Theorem~\ref{thm:mc2-bar-intrinsic}).
Averaging: $\operatorname{av}(D_\cA^{E_1}) = D_\cA$ because
the genus-completed bar differential on ${\Barch}^{\mathrm{ord}}(\cA)$
descends to $\Barch(\cA)$ under $\Sigma_n$-coinvariants, and
$\operatorname{av}(\dzero) = \dzero$.
\end{proof}

\begin{proposition}[\ClaimStatusProvedHere]
\label{prop:e1-shadow-r-matrix}
\index{shadow Postnikov tower!E1@$E_1$ variant}
\index{R-matrix!as $E_1$ shadow}
The arity-$r$ component of $\Theta_\cA^{E_1}$ at genus~$0$ is an
element $r_r \in
\operatorname{End}(V^{\otimes r}) \otimes H^{r-2}(K_r)
\otimes \cO(*\Delta)$, whose $\Sigma_r$-coinvariant is the
$E_\infty$ shadow at arity~$r$\textup{:}
\begin{center}
\renewcommand{\arraystretch}{1.15}
\begin{tabular}{lcc}
\toprule
\emph{Arity} & \emph{$E_1$ shadow} & \emph{$E_\infty$ shadow} \\
\midrule
$r = 2$ & $r(z)$ (classical $r$-matrix)
  & $\kappa(\cA)$ \\
$r = 3$ & $r_3(z_1,z_2)$ (associator shadow)
  & $\mathfrak{C}(\cA)$ \\
$r = 4$ & $r_4(z_1,z_2,z_3)$ (quartic $R$-matrix shadow)
  & $\mathfrak{Q}(\cA)$ \\
\bottomrule
\end{tabular}
\end{center}
The MC equation at arity~$r$ is the
$\partial K_{r+1}$-identity\textup{:}
CYBE \textup{(}$r{=}2$, two faces of $K_3$\textup{)},
pentagon \textup{(}$r{=}3$, five faces of $K_4$\textup{)},
quartic $R$-matrix identity \textup{(}$r{=}4$\textup{)}.
\end{proposition}

\begin{proof}
At genus~$0$, $\cM_\Ass(0,r) = C_*(K_r)$.  The MC equation
projected to arity~$r$ is the equation on
$\partial K_{r+1}$: each codimension-$1$ face of $K_{r+1}$
contributes one term.  The classical identifications
(Stasheff~\cite{Sta63}) give $|\partial K_3| = 2$ (CYBE),
$|\partial K_4| = 5$ (pentagon),
$|\partial K_5| = 9$ (quartic identity).  Coinvariant
projection: the averaging map sends the $\operatorname{End}(V^{\otimes r})$-valued
$r_r$ to the $\Sigma_r$-invariant scalar
$\Theta^{\leq r}_{0,r}$
(Corollary~\ref{cor:shadow-extraction}).
\end{proof}


%% ════════════════════════════════════════════════════════════════════════
%%  THE E₁ SHADOW POSTNIKOV TOWER
%% ════════════════════════════════════════════════════════════════════════

\section{The \texorpdfstring{$E_1$}{E1} shadow Postnikov tower}
\label{sec:e1-shadow-tower}
\index{shadow Postnikov tower!E1@$E_1$ variant|textbf}
\index{R-matrix!shadow tower}
\index{Drinfeld associator!as arity-3 E1 shadow@as arity-$3$ $E_1$ shadow}

\begin{definition}[$E_1$ shadow at arity~$r$]
\label{def:e1-shadow-arity-r}
\ClaimStatusProvedHere
\index{E1 shadow@$E_1$ shadow!arity $r$|textbf}
The \emph{$E_1$ shadow at arity~$r$} of a cyclic $E_1$-chiral
algebra $\cA$ is the component
\begin{equation}
\label{eq:e1-shadow-arity-r}
r_r(z_1, \ldots, z_{r-1})
\;:=\;
(\Theta_\cA^{E_1})_{0,r}
\;\in\;
\operatorname{End}(V^{\otimes r})
\otimes
\Omega^{r-2}\bigl(
\overline{\mathrm{FM}}_r^{\mathrm{ord}}(\mathbb{C})
\bigr),
\end{equation}
where $(\Theta_\cA^{E_1})_{0,r}$ is the genus-$0$, arity-$r$
component of the $E_1$ MC element
(Theorem~\ref{rem:e1-mc-element}).
The variables $z_1, \ldots, z_{r-1}$ are the ordered collision
coordinates on $\overline{\mathrm{FM}}_r^{\mathrm{ord}}(\mathbb{C})$
(one variable for each consecutive pair, after fixing
translation and scale).
\end{definition}

\begin{theorem}[$E_1$ MC equation at finite arity; \ClaimStatusProvedHere]
\label{thm:e1-mc-finite-arity}
\index{E1 shadow@$E_1$ shadow!MC equation}
The Maurer--Cartan equation for $\Theta_\cA^{E_1}$, projected to
genus~$0$ and arity~$r$, gives the following identities:
\begin{enumerate}[label=\textup{(\roman*)}]
\item \emph{Arity~$2$ \textup{(}$r = 2$\textup{)}:} The classical
  Yang--Baxter equation\textup{:}
  \begin{equation}
  \label{eq:e1-mc-arity2}
  [r_{12}(z_1), r_{13}(z_1 + z_2)]
  + [r_{12}(z_1), r_{23}(z_2)]
  + [r_{13}(z_1 + z_2), r_{23}(z_2)]
  \;=\; 0.
  \end{equation}
  This is the content of $d^2 = 0$ on
  $\overline{\mathrm{FM}}_3^{\mathrm{ord}}(\mathbb{C})$
  \textup{(}the Yang--Baxter equation from $d^2 = 0$).
\item \emph{Arity~$3$ \textup{(}$r = 3$\textup{)}:} The
  pentagon\slash quasi-associativity relation for the
  \emph{Drinfeld associator shadow}
  $r_3(z_1, z_2) \in \operatorname{End}(V^{\otimes 3})$\textup{:}
  \begin{equation}
  \label{eq:e1-mc-arity3}
  d_{\mathrm{bar}} r_3
  + [r_{12}, r_3]_{1,23}
  + [r_{23}, r_3]_{12,3}
  \;=\; 0,
  \end{equation}
  where the notation $[-, -]_{I,J}$ denotes the Lie bracket in
  which tensor factors in~$I$ are on the left and those in~$J$
  on the right.  The boundary faces of the associahedron $K_4$
  give the five pentagon terms.
\item \emph{Arity~$4$ \textup{(}$r = 4$\textup{)}:} The quartic
  $R$-matrix identity\textup{:} the obstruction to extending the
  triple $(r, r_3)$ to arity~$4$ is
  $o_5 \in H^2(F^5 {\gAmod}^{E_1} / F^6)$, and
  \begin{equation}
  \label{eq:e1-mc-arity4}
  d_{\mathrm{bar}} r_4
  + \bigl\{r_3 \circ_2 r_3 \text{ boundary terms}\bigr\}
  + [r, r_4]
  \;=\; 0.
  \end{equation}
\end{enumerate}
\end{theorem}

\begin{proof}
The MC equation
$D\Theta^{E_1} + \tfrac{1}{2}[\Theta^{E_1}, \Theta^{E_1}] = 0$
on ${\gAmod}^{E_1}$
(Definition~\ref{def:e1-modular-convolution})
projected to genus~$0$ and arity~$r$ uses the face structure of the
Stasheff associahedra $K_{r+1}$:
$\partial K_3$ has $2$ faces (CYBE);
$\partial K_4$ has $5$ faces (pentagon);
$\partial K_5$ has $9$ faces (the quartic identity).

Part~(i) is
the Yang--Baxter equation from $d^2 = 0$:
$K_3$ is the interval $[0,1]$, whose two endpoints give
the two terms of the CYBE.
Part~(ii): the genus-$0$, arity-$3$ MC component on
$\overline{\mathrm{FM}}_4^{\mathrm{ord}}(\mathbb{C})$ identifies
with the boundary analysis of $K_4$ (the pentagon), whose five
edges give the five terms of the quasi-associativity relation.
Part~(iii) follows by the same mechanism at arity~$4$,
using the nine codimension-$1$ faces of $K_5$.
\end{proof}

\begin{theorem}[Coinvariant projection: $E_1$ shadows to $E_\infty$ shadows; \ClaimStatusProvedHere]
\label{thm:e1-coinvariant-shadow}
\index{E1 shadow@$E_1$ shadow!coinvariant projection}
The $\Sigma_r$-coinvariant of the $E_1$ shadow at arity~$r$ is the
$E_\infty$ shadow at arity~$r$\textup{:}
\begin{align}
\label{eq:e1-coinvariant-arity2}
\operatorname{av}_{r=2}\bigl(r(z)\bigr)
&\;=\; \kappa(\cA), \\
\label{eq:e1-coinvariant-arity3}
\operatorname{av}_{r=3}\bigl(r_3(z_1, z_2)\bigr)
&\;=\; \mathfrak{C}(\cA), \\
\label{eq:e1-coinvariant-arity4}
\operatorname{av}_{r=4}\bigl(r_4(z_1, z_2, z_3)\bigr)
&\;=\; \mathfrak{Q}(\cA).
\end{align}
More precisely: the averaging map
$\operatorname{av}\colon {\gAmod}^{E_1} \twoheadrightarrow \gAmod$
of Definition~\textup{\ref{def:e1-modular-convolution}} sends each
$E_1$ shadow to the corresponding $E_\infty$ shadow at the same
arity.
\end{theorem}

\begin{proof}
The averaging map
$\operatorname{av}(\phi)(g,n) =
(1/n!)\sum_{\sigma \in \Sigma_n} \sigma \cdot \phi(g,n)$
sends $\operatorname{Hom}(\cM_\Ass(g,n), \operatorname{End}_\cA(n))$
to $\operatorname{Hom}_{\Sigma_n}(\cM_\Com(g,n),
\operatorname{End}_\cA(n))$.  Since the MC element
$\Theta^{E_1}_\cA$ maps to $\Theta_\cA$ under averaging
(eq.~\eqref{eq:e1-to-einfty-mc}), the component projections also
correspond: the arity-$r$ component of
$\operatorname{av}(\Theta^{E_1})$ is
$\operatorname{av}(r_r)$, and this equals $\Theta^{\leq r}_{0,r}$
by construction.  The identification of $\Theta^{\leq 2}_{0,2}$
with $\kappa(\cA)$, $\Theta^{\leq 3}_{0,3}$ with
$\mathfrak{C}(\cA)$, and $\Theta^{\leq 4}_{0,4}$ with
$\mathfrak{Q}(\cA)$ is
Definition~\ref{def:shadow-postnikov-tower}
  and Corollary~\ref{cor:shadow-extraction}.
\end{proof}

\begin{remark}[Finite-arity evidence for the FG shadow conjecture]
\label{rem:e1-coinvariant-shadow-fg}
Theorem~\ref{thm:e1-coinvariant-shadow} supplies the arity-$2,3,4$
verification of Conjecture~\ref{conj:FG-shadow}: at finite arity, the
Francis--Gaitsgory shadow is the $\Sigma_n$-coinvariant of the
associative theory.
\end{remark}

\begin{example}[The $E_1$ shadow archetypes]
\label{ex:e1-shadow-archetypes}
\index{E1 shadow@$E_1$ shadow!archetypes}
\begin{center}
\renewcommand{\arraystretch}{1.15}
\begin{tabular}{lccc}
\toprule
& \emph{$r(z)$}
& \emph{$r_3(z_1,z_2)$}
& \emph{$r_4(z_1,z_2,z_3)$} \\
\midrule
Heisenberg $\cH_k$
& $k/z$ (scalar)
& $0$
& $0$ \\
Affine $\hat\fg_k$
& $k\Omega/z$
& KZ associator $\Phi_{\mathrm{KZ}}(\hat\fg_k)$
& $0$ \\
$\beta\gamma$
& $0$
& $0$
& $\mathfrak{Q}^{\mathrm{ct}} \neq 0$ \\
Virasoro
& $c/(2z)$
& non-trivial
& non-trivial \\
$Y_\hbar(\fg)$
& $r_{\mathrm{YB}}(z)$
  (Yang--Baxter)
& Drinfeld associator $\Phi_{\hbar}(\fg)$
& non-trivial \\
\bottomrule
\end{tabular}
\end{center}
\end{example}

\begin{remark}[Three $r$-matrices: Laplace kernel, collision residue,
  pre-dualisation form]
\label{rem:three-r-matrices}
\index{r-matrix@$r$-matrix!disambiguation}
\index{r-matrix@$r$-matrix!Laplace kernel}
\index{r-matrix@$r$-matrix!collision residue}
\index{r-matrix@$r$-matrix!pre-dualisation form}
Three distinct objects share the notation~$r(z)$ in this text.
Each is a meromorphic function of the separation~$z = z_1 - z_2$;
they differ in pole orders, target spaces, and normalisation.
\begin{enumerate}[label=\textup{(\alph*)}]

\item \emph{Laplace kernel}
  $r^{\mathrm{L}}_\cA(z)
   = \int_0^\infty e^{-\lambda z}\,
     \{{\cdot}\,_\lambda\,{\cdot}\}\,d\lambda
   \in \cA\,\widehat\otimes\,\cA[\![z^{-1}]\!]$:
  the full OPE generating function, encoding every singular term of
  the operator product expansion as a Laurent series in~$z^{-1}$.
  Pole orders match the OPE directly.
  \begin{center}
  \renewcommand{\arraystretch}{1.1}
  \begin{tabular}{ll}
  Heisenberg $\cH_k$\textup: &
    $r^{\mathrm{L}}(z) = k/z^2$ \\
  Affine $\hat\fg_k$\textup: &
    $r^{\mathrm{L}}(z)
     = k\,\kappa^{ab}/z^2
       + {f^{ab}}_c\,J^c/z$ \\
  Virasoro\textup: &
    $r^{\mathrm{L}}(z)
     = (c/2)/z^4 + 2T/z^2 + \partial T/z$
  \end{tabular}
  \end{center}

\item \emph{Collision residue}
  $r^{\mathrm{coll}}_\cA(z)
   = \operatorname{Res}^{\mathrm{coll}}_{0,2}(\Theta_\cA)
   \in \cA^!\,\widehat\otimes\,\cA^!$:
  the bar-theoretic $r$-matrix, extracted by residue along the
  collision divisor using the $d\log(z_1-z_2)$ measure.
  Because $d\log$ absorbs one power of $(z_1-z_2)$, every pole order
  drops by one relative to the Laplace kernel.
  For bosonic algebras, only odd-order poles survive
  (even-order OPE poles descend to odd-order residues).
  The target space is $\cA^!\otimes\cA^!$: passage to
  the Koszul dual is part of the construction.
  \begin{center}
  \renewcommand{\arraystretch}{1.1}
  \begin{tabular}{ll}
  Heisenberg $\cH_k$\textup: &
    $r^{\mathrm{coll}}(z) = k/z$
    \quad\textup{(simple pole; OPE had double)} \\
  Affine $\hat\fg_k$\textup: &
    $r^{\mathrm{coll}}(z)
     = \Omega_\fg/((k{+}h^\vee)z)$
    \quad\textup{(Sugawara-normalised)} \\
  Virasoro\textup: &
    $r^{\mathrm{coll}}(z) = (c/2)/z^3 + 2T/z$
    \quad\textup{(cubic + simple; no even poles)}
  \end{tabular}
  \end{center}
  This is the master object from which the others descend:
  the genus-$0$, arity-$2$ shadow
  $r^{\mathrm{coll}}(z) = \mathrm{Sh}_{0,2}(\Theta_\cA)(z)$.

\item \emph{Pre-dualisation form}
  $r^{\mathrm{pre}}_\cA(z) \in \cA\otimes\cA$:
  the collision residue computed in $\cA\otimes\cA$, before
  applying the Koszul pairing to pass one or both factors to~$\cA^!$.
  Pole orders and pole structure are identical to~\textup{(b)};
  only the target space differs.
  \begin{center}
  \renewcommand{\arraystretch}{1.1}
  \begin{tabular}{ll}
  Heisenberg $\cH_k$\textup: &
    $r^{\mathrm{pre}}(z) = k/z$ \\
  Affine $\hat\fg_k$\textup: &
    $r^{\mathrm{pre}}(z) = k\,\Omega_\fg/z$ \\
  Virasoro\textup: &
    $r^{\mathrm{pre}}(z) = (c/2)/z^3 + 2T/z$
  \end{tabular}
  \end{center}
\end{enumerate}

\medskip
\noindent\textbf{Relations.}
The three are connected by two operations.
\begin{itemize}
\item $d\log$ extraction:
  $r^{\mathrm{L}}(z) \;\xmapsto{\;\operatorname{Res}_{d\log}\;}
   r^{\mathrm{pre}}(z)$.
  Every pole order drops by one.
  Mnemonic: $z^{-n}\,dz = z^{-(n-1)}\,d\log z$,
  so the residue along $d\log(z_1-z_2)$ of a pole of order~$n$
  in the OPE is a pole of order~$n-1$ in the collision residue.
\item Koszul dualisation:
  $r^{\mathrm{pre}}(z) \;\xmapsto{\;\mathrm{id}\otimes\kappa\;}
   r^{\mathrm{coll}}(z)$.
  Same pole structure, different target space
  ($\cA\otimes\cA \to \cA^!\otimes\cA^!$).
\end{itemize}
The $E_1$ scalar shadow
$r^{\mathrm{sc}}(z) \in \C(z)$ appearing in the archetype table
\textup{(}Example~\textup{\ref{ex:e1-shadow-archetypes})} is the
vacuum projection of $r^{\mathrm{coll}}(z)$; its
$\Sigma_2$-coinvariant is the scalar curvature
$\operatorname{av}(r^{\mathrm{sc}}(z)) = \kappa(\cA)$.

\medskip
\noindent\textbf{Normalisation note for affine Kac--Moody.}
The archetype table
\textup{(}Example~\textup{\ref{ex:e1-shadow-archetypes})}
lists $r(z) = k\Omega/z$: the pre-dualisation collision residue,
unnormalised.  Some formal constructions
\textup{(}e.g.\ the KZ connection,
eq.~\textup{\eqref{eq:kz-connection-bar})} use the
Sugawara-normalised form $\Omega/((k{+}h^\vee)z)$, which is
$r^{\mathrm{coll}}(z)$ after passage to~$\cA^!$.
Both are collision residues with a simple pole in~$z$; they
differ by the scalar factor $k/(k{+}h^\vee)$
encoding the Koszul-dual level shift $k \mapsto -k - 2h^\vee$.

\medskip
\noindent\textbf{Convention.}
When the text writes $r(z)$ unqualified, the context determines
which variant is meant.  In the $E_1$ shadow tower
\textup{(}this chapter\textup{)},
$r(z) = r^{\mathrm{sc}}(z)$ is the scalar shadow.
In the holographic datum
\textup{(}\S\textup{\ref{sec:thqg-gravitational-yangian})}
and Yangian constructions,
$r(z) = r^{\mathrm{coll}}(z)$ is the full collision residue.
In explicit OPE computations,
$r(z) = r^{\mathrm{L}}(z)$ is the Laplace kernel.
\end{remark}

\begin{remark}[Arithmetic refinement of $E_1$ shadow depth]
\label{rem:e1-arithmetic-shadow-depth}
The $E_1$ shadow depth $r_{\max}$ admits a finer decomposition
$d = 1 + d_{\mathrm{arith}} + d_{\mathrm{alg}}$ via
Theorem~\textup{\ref{thm:depth-decomposition}}
(Chapter~\textup{\ref{chap:arithmetic-shadows}}),
where $d_{\mathrm{arith}}$ counts critical lines of $L$-functions
governing the spectral curve
and $d_{\mathrm{alg}}$ records the algebraic shadow tower contribution.
\end{remark}

\begin{construction}[KZ associator as arity-$3$ $E_1$ shadow
of $\hat\fg_k$]
\label{constr:kz-associator-e1-shadow}
\ClaimStatusProvedHere
\index{KZ associator!as E1 shadow@as $E_1$ shadow}
\index{affine Kac--Moody!E1 shadow@$E_1$ shadow}
For $\cA = \hat\fg_k$ (affine Kac--Moody algebra), the ordered bar
complex on $\overline{\mathrm{FM}}_3^{\mathrm{ord}}(\mathbb{C})$
carries the KZ connection
\begin{equation}
\label{eq:kz-connection-bar}
\nabla_{\mathrm{KZ}}
\;=\;
d - \frac{1}{k+h^\vee}\,
\Bigl(
\frac{\Omega_{12}}{z_1-z_2}\,dz_{12}
+ \frac{\Omega_{23}}{z_2-z_3}\,dz_{23}
\Bigr),
\end{equation}
where $\Omega_{ij} = \sum_a t^a_i \otimes t^a_j$ is the Casimir
acting on the $i$-th and $j$-th tensor factors.  The monodromy of
this connection as the ordered pair $(z_1, z_2, z_3)$ traverses the
boundary cycle of
$\overline{\mathrm{FM}}_3^{\mathrm{ord}}(\mathbb{C})$ is the
\emph{KZ associator}
\begin{equation}
\label{eq:kz-associator-monodromy}
\Phi_{\mathrm{KZ}}
\;=\;
\sum_{w \in \{t_{12}, t_{23}\}^*}
\zeta(w)\, w
\;\in\;
\exp\bigl(\widehat{\mathrm{Lie}}(t_{12}, t_{23})\bigr),
\end{equation}
with $\zeta(w)$ the Drinfeld--Kontsevich multiple zeta values.
This is the arity-$3$ $E_1$ shadow:
\begin{equation}
\label{eq:r3-kz}
r_3(\hat\fg_k)
\;=\;
\Phi_{\mathrm{KZ}}(\hat\fg_k)
\;=\;
\Phi_{\mathrm{KZ}}\!\bigl(
\tfrac{1}{k+h^\vee}\, \Omega_{12},\;
\tfrac{1}{k+h^\vee}\, \Omega_{23}
\bigr).
\end{equation}
The coinvariant $\operatorname{av}(\Phi_{\mathrm{KZ}}) =
\mathfrak{C}(\hat\fg_k)$ is the cubic shadow: the
$\Sigma_3$-averaged monodromy reduces to
$\kappa(\hat\fg_k) \cdot (x,[y,z])$.
\end{construction}

\begin{construction}[Formal ordered arity-$2$ shadow series]
\label{constr:modular-r-matrix-genus1}
\ClaimStatusProvedHere
\index{arity-2 shadow!ordered genus refinement}
\index{modular R-matrix@modular $R$-matrix!ordered shadow surface}
The genus-refined arity-$2$ projection of the ordered $E_1$
Maurer--Cartan element defines the formal series
\begin{equation}
\label{eq:modular-r-matrix-genus1}
R^{E_1,\mathrm{bin}}(z; \hbar)
\;=\;
\sum_{g \geq 0} \hbar^{2g}\, r_g(z)
\end{equation}
with genus-$0$ term $r_0(z)=R(z)$ equal to the ordered genus-$0$
exchange kernel.  This series is the formal ordered arity-$2$ shadow
package attached to $\Theta^{E_1}_\cA$.

Its coefficientwise Maurer--Cartan relations come from the arity-$2$
projection of the ordered $E_1$ MC equation.  On affine or free
benchmark surfaces one expects elliptic/genus-$1$ formulas for the
first correction $r_1(z)$, but such line-side modular-kernel
realizations are not part of the proved theorem surface of this
chapter.
\end{construction}


%% ════════════════════════════════════════════════════════════════════════
%%  THE E₁ FIVE MAIN THEOREMS AT ALL GENERA
%% ════════════════════════════════════════════════════════════════════════

\section{The \texorpdfstring{$E_1$}{E1} five main theorems at
all genera}
\label{sec:e1-five-theorems}
\index{five main theorems!E1@$E_1$ variant!all genera|textbf}

\begin{theorem}[Theorem~$\mathrm{A}^{E_1}$ at all genera: ordered bar--cobar adjunction; \ClaimStatusProvedHere]
\label{thm:e1-theorem-A-modular}
\label{thm:e1-theorem-A}
\index{five main theorems!E1@$E_1$ variant!Theorem A (modular)}
Let $\cA$ be a cyclic $E_1$-chiral algebra.  The genus-$g$ ordered
bar complex ${\Barch}^{\mathrm{ord}}(\cA)(g,n)$ is an
$F\!\Ass$-coalgebra
\textup{(}Definition~\textup{\ref{def:feynman-transform-ass}})
with differential satisfying $D^2 = 0$.
The bar and cobar functors at genus~$g$ form an adjunction
\[
\Cobar^{(g)}
\;\dashv\;
{\Barch}^{(g)}
\;\colon\;
\mathsf{Alg}_{F\!\Ass}
\;\rightleftarrows\;
\mathsf{Coalg}_{F\!\Ass}.
\]
Koszul exchange carries the formal ordered arity-$2$ shadow series
$R^{E_1,\mathrm{bin}}(z;\hbar)$
\textup{(}Construction~\textup{\ref{constr:modular-r-matrix-genus1}}):
its genus-$g$ coefficient records the ordered arity-$2$ correction on
this bar-level surface.  Interpreting this series as a genuine
higher-genus line-side modular $R$-matrix requires extra Yangian input
not proved here.
\end{theorem}

\begin{proof}
The $F\!\Ass$-coalgebra structure on the ordered bar complex
follows from the universal property of the Feynman transform:
$D_{F\!\Ass}^2 = 0$
(Theorem~\ref{thm:fass-d-squared-zero}) gives the differential
on the graph complex, and the bar differential of $\cA$ at each
vertex is compatible with the operad composition maps of
$\cM_\Ass$.  The adjunction follows by the same argument as in
the $E_\infty$ case
(Theorem~\ref{thm:bar-cobar-adjunction}),
replacing the commutative modular operad by $\cM_\Ass$:
the cobar construction $\Cobar^{(g)}$ is defined as the tensor
algebra over $F\!\Ass$ with differential dual to the coalgebra
structure.

The formal ordered arity-$2$ shadow package: at each genus~$g$, the
opposite-duality identification
${\Barch}^{\mathrm{ord}}(\cA^{\op}) \cong
{\Barch}^{\mathrm{ord}}(\cA)^{\cop}$
is encoded at the arity-$2$ shadow level by
$R^{E_1,\mathrm{bin}}(z;\hbar)$ restricted to genus
$\leq g$.  This is the $E_1$ analogue of Verdier intertwining on the
ordered bar surface.
\end{proof}

\begin{theorem}[Theorem~$\mathrm{B}^{E_1}$ at all genera: ordered bar--cobar inversion; \ClaimStatusProvedHere]
\label{thm:e1-theorem-B-modular}
\label{thm:e1-theorem-B}
\index{five main theorems!E1@$E_1$ variant!Theorem B (modular)}
On the $E_1$ Koszul locus, the counit
\[
\Cobar^{(g)}\bigl({\Barch}^{\mathrm{ord}}(\cA)(g,\bullet)\bigr)
\;\xrightarrow{\;\sim\;}\;
\cA
\]
is a quasi-isomorphism at every genus~$g$, and the
$E_1$ Koszul dual is an $E_1$-chiral algebra at all genera.
\end{theorem}

\begin{proof}
PBW-completeness transfers from genus~$0$ to genus~$g$ via the
MC perturbation lemma: the genus-$g$ differential is a perturbation
of the genus-$0$ differential by curvature $\kappa \cdot \omega_g$,
and this perturbation preserves the PBW filtration (the curvature is
in the center of the algebra).  The inversion theorem for perturbed
complexes (the filtered quasi-isomorphism criterion of
Theorem~\ref{thm:bar-cobar-inversion-qi}) then gives the
result at each genus.
\end{proof}

\begin{theorem}[Theorem~$\mathrm{C}^{E_1}$ at all genera: ordered complementarity; \ClaimStatusProvedHere]
\label{thm:e1-theorem-C-modular}
\label{thm:e1-theorem-C}
\index{five main theorems!E1@$E_1$ variant!Theorem C (modular)}
At genus~$g$, the ordered complementarity takes the form
\begin{equation}
\label{eq:e1-complementarity-modular}
Q_g^{E_1}(\cA) + Q_g^{E_1}(\cA^{!,E_1})
\;=\;
H^*\!\bigl(\overline{\cM}_{g,n}^{\,\mathrm{rib}},\,
Z^{E_1}(\cA)\bigr),
\end{equation}
where $Q_g^{E_1}(\cA) :=
C_\Delta^{(g)} \square_{C^{e,(g)}}^{\mathbf{R}} C_\Delta^{(g)}$
is the genus-$g$ ordered partition function and
$Z^{E_1}(\cA)$ is the $E_1$ center \textup{(}the space of
$R$-matrix-equivariant elements\textup{)}.
\end{theorem}

\begin{proof}
The genus-$g$ ordered complementarity extends from genus~$0$
(the genus-$0$ ordered complementarity theorem) by the same modular extension
principle used for the $E_\infty$ complementarity
(Theorem~\ref{thm:quantum-complementarity-main}): the
$E_1$ MC element $\Theta^{E_1}_\cA$ gives a genus-$g$
deformation of the genus-$0$ diagonal bicomodule, and the
ordered Hochschild--coHochschild identification carries through
at each genus.  The $E_1$ center $Z^{E_1}(\cA)$ is the subspace
of elements satisfying $R(z) \cdot a = a$ (braiding-invariant),
replacing the naive center of the $E_\infty$ theory.
\end{proof}

\begin{theorem}[Theorem~$\mathrm{D}^{E_1}$ at all genera: formal ordered arity-$2$ shadow series; \ClaimStatusProvedHere]
\label{thm:e1-theorem-D-modular}
\label{thm:e1-theorem-D}
\index{five main theorems!E1@$E_1$ variant!Theorem D (modular)}
\index{arity-2 shadow!as $E_1$ characteristic package}
The formal ordered arity-$2$ shadow series
$R^{E_1,\mathrm{bin}}(z;\hbar) = \sum_{g \geq 0}
\hbar^{2g}\,r_g(z)$
obtained from the genus-refined arity-$2$ projection of
$\Theta^{E_1}_\cA$ is the $E_1$ arity-$2$ characteristic package of
$\cA$.  It satisfies the arity-$2$ projection of the ordered
Maurer--Cartan equation
\begin{equation}
\label{eq:modular-ybe}
D\,R^{E_1,\mathrm{bin}} +
\tfrac{1}{2}[R^{E_1,\mathrm{bin}}, R^{E_1,\mathrm{bin}}]
\;=\; 0
\end{equation}
in ${\gAmod}^{E_1}$, and its properties mirror those of the scalar
$\kappa(\cA)$\textup{:}
\begin{enumerate}[label=\textup{(\roman*)}]
\item \emph{Universality}\textup{:}
  $R^{E_1,\mathrm{bin}}$ depends only on the $E_1$-chiral algebra
  $\cA$, not on auxiliary choices.
\item \emph{Additivity}\textup{:} for $\cA = \cA_1 \otimes
  \cA_2$ with vanishing mixed OPE,
  $R^{E_1,\mathrm{bin}}_{\cA_1 \otimes \cA_2}
  = R^{E_1,\mathrm{bin}}_{\cA_1} \otimes 1 + 1 \otimes
  R^{E_1,\mathrm{bin}}_{\cA_2}$.
\item \emph{Anti-symmetry}\textup{:}
  $R^{E_1,\mathrm{bin}}_{\cA^{\op}}(z;\hbar)
  = R^{E_1,\mathrm{bin}}_{\cA}(z;\hbar)^{-1}$ on the ordered
  bar surface.
\item \emph{Coinvariant recovery}\textup{:}
  $\operatorname{av}\bigl(R^{E_1,\mathrm{bin}}(z;\hbar)\bigr)
  = \kappa(\cA) + O(\hbar^2)$.
\end{enumerate}
\noindent
This theorem packages the ordered arity-$2$ shadow data only.  A
genuine higher-genus line-side modular $R$-matrix interpretation is
additional Yangian input and is not proved here.
\end{theorem}

\begin{proof}
The arity-$2$ ordered Maurer--Cartan relation~\eqref{eq:modular-ybe} is the
arity-$2$ MC equation in ${\gAmod}^{E_1}$: it follows from
$D_{F\!\Ass}^2 = 0$ by restriction to the two-point
component.

(i)~Universality: $R^{E_1,\mathrm{bin}}$ is determined by the
$F\!\Ass$-coalgebra structure on ${\Barch}^{\mathrm{ord}}(\cA)$,
which depends only on~$\cA$.

(ii)~Additivity: for $\cA = \cA_1 \otimes \cA_2$ with vanishing
mixed OPE, the ordered bar complex decomposes as a product
${\Barch}^{\mathrm{ord}}(\cA_1 \otimes \cA_2) \cong
{\Barch}^{\mathrm{ord}}(\cA_1) \otimes
{\Barch}^{\mathrm{ord}}(\cA_2)$
(ordered shuffle theorem), and the monodromy of a product
connection is additive.

(iii)~Anti-symmetry: opposite-duality gives
${\Barch}^{\mathrm{ord}}(\cA^{\op}) \cong
{\Barch}^{\mathrm{ord}}(\cA)^{\cop}$,
and the monodromy of the co-opposite is the inverse.

(iv)~Coinvariant: by
Theorem~\ref{thm:e1-coinvariant-shadow} at arity~$2$.
\end{proof}

\begin{theorem}[Theorem~$\mathrm{H}^{E_1}$ at all genera: ordered Hochschild at genus~$g$; \ClaimStatusProvedHere]
\label{thm:e1-theorem-H-modular}
\label{thm:e1-theorem-H}
\index{five main theorems!E1@$E_1$ variant!Theorem H (modular)}
For every genus~$g$ and complete $\cA$-bimodule~$M$, the
genus-$g$ ordered Hochschild homology is computed
coalgebraically\textup{:}
\[
\HH^{\mathrm{ch},(g)}_\bullet(\cA, M)
\;\simeq\;
\coHoch^{\mathrm{ch},(g)}_\bullet\!\bigl(
C, \KK_\cA^{\mathrm{bi}}(M)
\bigr),
\]
where $C = {\Barch}^{\mathrm{ord}}(\cA)$ and the genus-$g$ twisted
differential uses the curvature $\kappa(\cA) \cdot \omega_g$.
\end{theorem}

\begin{proof}
The genus-$g$ extension of the Hochschild--coHochschild dictionary
(the genus-$0$ ordered Hochschild theorem) follows by the same modular
perturbation argument as
Theorem~\ref{thm:e1-theorem-B-modular}: the genus-$g$
perturbation preserves the bimodule-bicomodule equivalence.
\end{proof}

\begin{remark}[Resolution of the ordered associative modular
Maurer--Cartan theorem]
\label{rem:conj-modular-resolved}
\index{conj:modular!resolved}
Theorems~\textup{\ref{thm:e1-theorem-A-modular}}--\textup{%
\ref{thm:e1-theorem-H-modular}} establish the associative modular
MC theory asked for by
Theorem~\textup{\ref{thm:ordered-associative-modular-mc}}: the $E_1$ deformation complex
is ${\gAmod}^{E_1}$
\textup{(}Definition~\textup{\ref{def:e1-modular-convolution}}),
and the canonical MC element is $\Theta_\cA^{E_1} := D_\cA^{E_1}
- \dzero$.  The ``ordered cyclic compactification'' sought in the
proof-obligations remark is exactly the Feynman transform $F\!\Ass$
of the ribbon modular operad
\textup{(}Definition~\textup{\ref{def:feynman-transform-ass}}).
\end{remark}

\begin{remark}[$E_1$ primitive kernel]
\label{rem:e1-primitive-kernel}
\index{primitive kernel!E1@$E_1$ variant}
\index{E1 algebra@$E_1$-algebra!primitive kernel}
The cofree-coderivation principle for $F\!\Ass$ yields an
\emph{$E_1$ primitive kernel}
\begin{equation}
\label{eq:e1-primitive-kernel}
K_\cA^{E_1}
\;:=\;
\log_\star\bigl(\Theta_\cA^{E_1}\bigr)
\;\in\;
\prod_{\substack{g,n \\ 2g-2+n > 0}}
\operatorname{End}(V^{\otimes n})
\otimes
\Omega^{n-2}\!\bigl(
\overline{\mathrm{FM}}_n^{\,\mathrm{ord}}(\mathbb{C})
\bigr),
\end{equation}
where $\log_\star$ is the operadic logarithm on the
$F\!\Ass$-coalgebra ${\Barch}^{\mathrm{ord}}(\cA)$, decomposing the
full MC element into primitive (one-vertex) contributions.

\medskip
\noindent\textbf{Coinvariant recovery.}
The averaging map of
Definition~\textup{\ref{def:e1-modular-convolution}} sends the $E_1$
primitive kernel to the $E_\infty$ primitive kernel:
\begin{equation}
\label{eq:e1-coinvariant-kernel}
\operatorname{av}\bigl(K_\cA^{E_1}\bigr)
\;=\;
K_\cA.
\end{equation}
This is the primitive-level refinement of
Theorem~\textup{\ref{thm:e1-coinvariant-shadow}}: the coinvariant
projection is compatible with the operadic logarithm because
$\operatorname{av}$ is a dg~Lie morphism.

\medskip
\noindent\textbf{Low-arity identifications.}
\begin{enumerate}[label=\textup{(\roman*)}]
\item $(K_\cA^{E_1})_{0,2} = r(z)$, the classical $r$-matrix
  \textup{(}genus~$0$, arity~$2$\textup{)}.  The coinvariant
  $\operatorname{av}(r(z)) = \kappa(\cA)$ recovers the scalar
  curvature.
\item $(K_\cA^{E_1})_{0,3} = \Phi_{\mathrm{KZ}}(\cA)$, the
  KZ\slash Drinfeld associator \textup{(}genus~$0$, arity~$3$,
  for affine algebras;
  Construction~\textup{\ref{constr:kz-associator-e1-shadow}}).
  Coinvariant: $\operatorname{av}(\Phi_{\mathrm{KZ}}) =
  \mathfrak{C}(\cA)$, the cubic shadow.
\item $(K_\cA^{E_1})_{1,1}$ is the genus-$1$ primitive, whose
  coinvariant gives the genus-$1$ curvature $\kappa(\cA) \cdot
  \lambda_1$.
\end{enumerate}

\medskip
\noindent\textbf{$E_1$ primitive master equation.}
The MC equation $D\Theta^{E_1} + \tfrac{1}{2}[\Theta^{E_1},
\Theta^{E_1}] = 0$ expressed in primitive generators takes the form
\begin{equation}
\label{eq:e1-primitive-master-equation}
dK_\cA^{E_1}
\;+\;
K_\cA^{E_1} \star_{E_1} K_\cA^{E_1}
\;=\; 0,
\end{equation}
where $\star_{E_1}$ is the convolution product on the $F\!\Ass$
modular cooperad, whose arity-$r$ component is governed by the face
structure of the Stasheff associahedron $K_{r+1}$.  In particular,
the arity-$2$ component of~\eqref{eq:e1-primitive-master-equation} is
the classical Yang--Baxter equation, and the arity-$3$ component is
the pentagon equation.  Averaging
recovers the $E_\infty$ primitive master equation
$dK_\cA + K_\cA \star K_\cA
+ \hbar\,\Delta_{\mathrm{ns}}(K_\cA) = 0$.
\end{remark}

\section{Cyclic trace, ribbon structure, and the
\texorpdfstring{$1/N$}{1/N} expansion}%
\label{sec:e1-ribbon-thooft}%
\index{ribbon structure!from cyclic trace|textbf}%
\index{t Hooft expansion@'t~Hooft expansion!from ribbon structure|textbf}

\begin{lemma}[Bare graphs do not determine a 't~Hooft expansion; \ClaimStatusProvedHere]%
\label{lem:bare-graph-no-thooft}%
\index{t Hooft expansion@'t~Hooft expansion!bare graph obstruction}%
A bare modular stable-graph expansion does not
canonically define a 't~Hooft $N$-expansion.
\end{lemma}

\begin{proof}
The $N$-weight in a 't~Hooft expansion is
$N^{V - E + F} = N^{\chi(\Sigma_\Gamma)}$,
determined by the number of faces~$F$ of the
thickened surface~$\Sigma_\Gamma$.
A bare stable graph does not specify cyclic orderings
of half-edges at vertices; hence it does not determine
a ribbon surface, hence does not determine~$F$,
hence carries no canonical $N$-power.
The same underlying graph admits inequivalent ribbon
structures with different genera.
\end{proof}

\begin{theorem}[Cyclicity is the ribbon-enabling datum;
\ClaimStatusProvedHere]%
\label{thm:cyclicity-ribbon}%
\index{cyclic A-infinity@cyclic $A_\infty$!ribbon structure|textbf}%
\index{ribbon structure!from cyclicity|textbf}%
Let $(O, m_\bullet, \tau)$ be a graded vector space
with multilinear operations
$m_n\colon O^{\otimes n} \to O$ and a
degree-zero nondegenerate bilinear pairing
$\tau\colon O \otimes O \to k$.
Define the cyclic tensors
\begin{equation}\label{eq:cyclic-tensors}
\tilde{m}_n(x_0,\ldots,x_n)
\;:=\;
\tau\!\bigl(m_n(x_1,\ldots,x_n),\,x_0\bigr).
\end{equation}
If $\tilde{m}_n$ satisfies cyclic invariance
$\tilde{m}_n(x_0,\ldots,x_n) = (-1)^\epsilon\,
\tilde{m}_n(x_1,\ldots,x_n,x_0)$
\textup{(}with Koszul sign~$\epsilon$\textup{)}, then
each operation~$m_n$ determines a cyclic
$(n{+}1)$-tensor, and every Feynman graph built from
these tensors canonically thickens to a ribbon graph.
\end{theorem}

\begin{proof}
The pairing identifies
$\Hom(O^{\otimes n}, O) \cong
\Hom(O^{\otimes(n+1)}, k)$.
Under this identification, $m_n$ becomes the
$(n{+}1)$-linear form~$\tilde{m}_n$.
Cyclic invariance means $\tilde{m}_n$ descends to a
function on cyclic words of length $n{+}1$:
there is no distinguished output among the $n{+}1$ legs.
Each vertex therefore carries a cyclic order on its
incident half-edges, which is exactly the data needed
for canonical ribbon thickening: at each vertex, the
cyclic order determines the local surface, and gluing
via internal edges produces a global ribbon surface
$\Sigma_\Gamma$ with a well-defined face decomposition.
\end{proof}

\begin{corollary}[Operads are too small for traces;
\ClaimStatusProvedHere]%
\label{cor:operads-too-small}%
\index{wheeled operad!necessity of}%
Any open/closed theory closed under Swiss-cheese
composition \emph{and} trace-contraction cannot be
encoded by a rooted operad.  The minimal closure is
a wheeled two-colored operadic object\textup;
for multi-interface gluing, a properad.
\end{corollary}

\begin{proof}
Operadic composition is indexed by rooted trees
\textup{(}acyclic, one output per vertex\textup{)}.
Trace-contraction $\sum_i m_2(x, e_i) \otimes e^i$
creates a one-vertex graph with a loop, not an operadic
graph.  The closure under both composition and
contraction generates graphs with arbitrary numbers of
loops: wheeled graphs.  Properadic structure allows
multiple outputs per vertex, accommodating
multi-boundary and multi-interface gluing.
\end{proof}

\begin{theorem}[Exact $N^{\chi}$ weighting from traced open color; \ClaimStatusProvedHere]%
\label{thm:exact-n-chi-weighting}%
\index{t Hooft expansion@'t~Hooft expansion!exact weighting|textbf}%
Let $(O, m_\bullet, \tau)$ be a cyclic
$A_\infty$-algebra with matrix realization
$O_N := \operatorname{Mat}_N \otimes O_0$ and trace
$\tau_N := \operatorname{Tr}_N \otimes \tau_0$.
For a connected ribbon graph~$\Gamma$ with
$V$~vertices, $E$~edges, and $F$~faces, the bare
amplitude scales as~$N^F$ \textup(one factor per
closed index loop\textup).
Under 't~Hooft rescaling
$g^2 \mapsto \lambda/N$, the rescaled amplitude
factors as
\begin{equation}\label{eq:n-chi-amplitude}
\operatorname{Amp}_\Gamma^{\lambda}(O_N)
\;=\;
N^{V - E + F}\cdot
\operatorname{Amp}_\Gamma^{\mathrm{red}}(O_0,\lambda)
\;=\;
N^{\chi(\Sigma_\Gamma)}\cdot
\operatorname{Amp}_\Gamma^{\mathrm{red}}(O_0,\lambda).
\end{equation}
The total genus expansion is then
\begin{equation}\label{eq:thooft-literal}
\sum_{\Gamma\;\text{ribbon}}
\frac{N^{\chi(\Sigma_\Gamma)}}{|\operatorname{Aut}_{\mathrm{rib}}(\Gamma)|}
\,W_\Gamma\,\Phi_\Gamma^T
\;=\;
\sum_{g=0}^\infty N^{2-2g}\,F_g(\lambda),
\end{equation}
where $\lambda = g_{\mathrm{YM}}^2 N$.  This is a
literal 't~Hooft expansion\textup: the $N$-power
is determined by the Euler characteristic of the
thickened surface.
\end{theorem}

\begin{proof}
Cyclic structure
\textup{(}Theorem~\textup{\ref{thm:cyclicity-ribbon}}\textup{)}
gives canonical ribbon thickening, so faces are
well-defined.  We count $N$-powers in two steps.

\emph{Step~1: bare double-line counting.}  With the
normalization $\tau_N = \operatorname{Tr}_N$ (not
$(1/N)\operatorname{Tr}_N$), the propagator
$\tau_N^{-1} = \sum_{i,j} E_{ij}\otimes E_{ji}$
carries no explicit $N$-factor: it merely routes
double-line indices through each internal edge.
Likewise, vertex traces route indices but produce no
bare $N$-powers.  The only source of $N$ is the
summation over each closed index loop: every face of
the ribbon graph~$\Gamma$ corresponds to one such
loop, contributing a factor of~$N$.  Hence the bare
amplitude scales as $N^F$.

\emph{Step~2: 't~Hooft rescaling.}  Writing the action
with overall normalization $g^{-2}\tau_N(\cdots)$,
each propagator acquires $g^2$ and each vertex
$g^{-2}$, giving total coupling $g^{2(E-V)}$.
Setting $\lambda = g^2 N$ yields
$g^{2(E-V)} = (\lambda/N)^{E-V}$, so the rescaled
amplitude carries $N$-power $F - (E - V) = V - E + F
= \chi(\Sigma_\Gamma)$ by Euler's formula.
\end{proof}

\begin{remark}[The logical chain]%
\label{rem:thooft-logical-chain}%
\index{t Hooft expansion@'t~Hooft expansion!logical chain}%
\begin{equation}\label{eq:thooft-chain}
\text{modular graph sum}
\;+\;
\text{cyclic trace}
\;+\;
\text{matrix realization}
\;\Longrightarrow\;
\text{literal 't~Hooft sum.}
\end{equation}
Plain modular $E_1$ gives genus-organized expansion.
Cyclic traced modular $E_1$ gives 't~Hooft expansion.
Ribbonized modular Swiss-cheese gives holographic
open/closed 't~Hooft language.
\end{remark}

\begin{definition}[Ribbonized modular Swiss-cheese theory]%
\label{def:ribbonized-swiss-cheese}%
\index{ribbonized Swiss-cheese|textbf}%
\index{Swiss-cheese!ribbonized modular}%
The \emph{ribbonized modular Swiss-cheese theory}
$\mathrm{RSC}_{\mathrm{tr}}^{\mathrm{mod}}$
is the two-colored wheeled modular properad with\textup:
\begin{enumerate}[label=\textup{(\roman*)}]
\item \emph{closed color}: the modular/chiral sector
  \textup{(}Chapters~\textup{\ref{sec:koszul-across-genera}--\ref{chap:e1-modular-koszul}}\textup{)}\textup;
\item \emph{open color}: the $E_1$/chiral or
  $A_\infty$-chiral line sector\textup;
\item \emph{mixed operations}: bulk-to-boundary and
  boundary-to-line interactions\textup;
\item \emph{trace data}: a nondegenerate invariant
  pairing~$\tau$ on the open color\textup;
\item \emph{wheel operations}: contractions of
  outputs with inputs via~$\tau$\textup;
\item \emph{ribbon data}: cyclic order at each open
  vertex, so every open graph thickens to a surface
  \textup{(}Theorem~\textup{\ref{thm:cyclicity-ribbon}}\textup{)}.
\end{enumerate}
The open input is a cyclic $A_\infty$-object
$(O, m_\bullet, \tau)$\textup; the closed input is
the modular datum $\Theta_\cA$.
The tree part recovers the Boardman--Vogt resolution
$W(\mathrm{SC}^{\mathrm{ch,top}})$ of the Swiss-cheese
operad \textup{(}Volume~II, \S\ref{cor:clean-replacement}\textup{)}.
\end{definition}

\begin{remark}[The correct holographic language]%
\label{rem:correct-holographic-language}%
\index{holographic language!correct form}%
The three levels of language are\textup:
\begin{center}
\small
\renewcommand{\arraystretch}{1.3}
\begin{tabular}{ll}
\toprule
\textbf{Language} & \textbf{Content} \\
\midrule
plain modular $E_1$
  & genus-organized expansion \\
cyclic traced modular $E_1$
  & 't~Hooft $1/N$ expansion \\
ribbonized modular Swiss-cheese
  & holographic open/closed 't~Hooft language \\
\bottomrule
\end{tabular}
\end{center}
\end{remark}
